% ---------- ABSTRACT (spec section 4) : page iii, <= 500 words ----------
\phantomsection\addcontentsline{toc}{chapter}{\protect\hspace*{0.5in}Abstract}
\frontheading{Abstract}

Locating people in aerial imagery is central to disaster search and rescue, but
the targets are often only a few pixels wide, partially occluded, and set against
cluttered backgrounds, which defeats standard object detectors. This thesis asks
which changes to a modern detector actually help for this task, and answers in
two phases: a controlled architecture ablation, and an extended study that
hardens the protocol and carries the model to real footage. Starting from a
YOLO11m baseline, an additive ablation introduces one change at a time and
evaluates every configuration on the C2A disaster-imagery dataset under one
training and evaluation protocol. The changes are channel-and-spatial attention
(CBAM) in place of the baseline self-attention block, a high-resolution P2
detection scale at stride four, and a bidirectional state-space (Mamba) neck.
The P2 scale is the effective change: it raises recall of the smallest
($<8$~px) targets from $0.743$ to $0.757$ and $\mathrm{AP}_{50}$ from $0.843$
to $0.853$ at a cost of about half a million parameters. CBAM is nearly free
and lowers the parameter count below the baseline while holding accuracy. The
state-space neck adds $2.4$~million parameters and roughly triples inference
latency without improving any accuracy metric, and is excluded. The recommended
CBAM~+~P2 detector, at $19.6$~million parameters, matches or exceeds all but one
published detector on the same dataset at about a third of its size, and runs in
real time on a single consumer GPU. The extended study then re-examines the
ground these results stand on. An audit finds that the official split shares
98\% of its test scenes with training; a scene-disjoint re-split costs $1.6$
points of $\mathrm{AP}_{50}$ and becomes the protocol for all further work.
Under it, the productive lever moves from architecture to training supervision:
blending a normalised Wasserstein similarity into the task-aligned label
assigner raises sub-8-pixel recall by $2.1$ points (95\% confidence interval
$[1.8, 2.4]$) at a measured cost on larger sizes, while a frequency-gated
evidence branch fails the same paired test and is reported as a second negative
result beside the state-space neck. Real aerial footage collected and annotated
for this thesis (drone flights at three altitudes, occluded campus scenes, and
frames from news video of disaster events) then measures transfer. Zero-shot,
the synthetic-trained model reaches $\mathrm{AP}_{50}$ $0.335$ on drone footage
and $0.133$ on real disaster frames; supervised fine-tuning on a few hundred
labelled real frames raises these to $0.795$ and $0.411$ while the source
benchmark moves by less than a point. A disaster event held out from all
training bounds the claim: there the adapted model's gain is not statistically
significant, and sub-8-pixel recall on real disaster footage is zero for every
model tested. The thesis contributes a controlled account of where small-object
gains come from on this task, a leakage-audited protocol, an assignment-level
mechanism with quantified trade-offs, an annotated real-imagery test suite, and
honest bounds on what a small labelling budget buys. The same lesson appears
twice: training supervision moved the metric where added architecture did not.
